% Options for packages loaded elsewhere
\PassOptionsToPackage{unicode}{hyperref}
\PassOptionsToPackage{hyphens}{url}
\PassOptionsToPackage{dvipsnames,svgnames,x11names}{xcolor}
\documentclass[
  10pt,
  letterpaper,
]{article}
\usepackage{xcolor}
\usepackage[margin=0.85in]{geometry}
\usepackage{amsmath,amssymb}
\setcounter{secnumdepth}{5}
\usepackage{iftex}
\ifPDFTeX
  \usepackage[T1]{fontenc}
  \usepackage[utf8]{inputenc}
  \usepackage{textcomp} % provide euro and other symbols
\else % if luatex or xetex
  \usepackage{unicode-math} % this also loads fontspec
  \defaultfontfeatures{Scale=MatchLowercase}
  \defaultfontfeatures[\rmfamily]{Ligatures=TeX,Scale=1}
\fi
\usepackage{lmodern}
\ifPDFTeX\else
  % xetex/luatex font selection
\fi
% Use upquote if available, for straight quotes in verbatim environments
\IfFileExists{upquote.sty}{\usepackage{upquote}}{}
\IfFileExists{microtype.sty}{% use microtype if available
  \usepackage[]{microtype}
  \UseMicrotypeSet[protrusion]{basicmath} % disable protrusion for tt fonts
}{}
\makeatletter
\@ifundefined{KOMAClassName}{% if non-KOMA class
  \IfFileExists{parskip.sty}{%
    \usepackage{parskip}
  }{% else
    \setlength{\parindent}{0pt}
    \setlength{\parskip}{6pt plus 2pt minus 1pt}}
}{% if KOMA class
  \KOMAoptions{parskip=half}}
\makeatother
\setlength{\parindent}{1em}
\setlength{\parskip}{2pt plus 1pt minus 1pt}
\usepackage{longtable,booktabs,array,tabularx}
\usepackage{caption}
\captionsetup[table]{skip=6pt}
\newcounter{none} % for unnumbered tables
\usepackage{calc} % for calculating minipage widths
% Correct order of tables after \paragraph or \subparagraph
\usepackage{etoolbox}
\makeatletter
\patchcmd\longtable{\par}{\if@noskipsec\mbox{}\fi\par}{}{}
\makeatother
% Allow footnotes in longtable head/foot
\IfFileExists{footnotehyper.sty}{\usepackage{footnotehyper}}{\usepackage{footnote}}
\makesavenoteenv{longtable}
\setlength{\emergencystretch}{3em} % prevent overfull lines
\providecommand{\tightlist}{%
  \setlength{\itemsep}{0pt}\setlength{\parskip}{0pt}}
\usepackage{bookmark}
\IfFileExists{xurl.sty}{\usepackage{xurl}}{} % add URL line breaks if available
\urlstyle{same}
% fallback for those not using the hyperref driver hyperxmp:
\makeatletter
\@ifundefined{xmpquote}{\newcommand{\xmpquote}[1]{#1}}{}
\makeatother
\hypersetup{
  pdftitle={Memory State: Inference-Native Retrieval for Long-Term Agent Memory},
  pdfauthor={Gordon Xiong},
  colorlinks=true,
  linkcolor={blue},
  filecolor={Maroon},
  citecolor={Blue},
  urlcolor={blue},
  pdfcreator={LaTeX via pandoc}}

\title{Memory State: Inference-Native Retrieval for Long-Term Agent
Memory}
\author{Gordon Xiong\\\small Seetie Inc.}
\date{\small Code and results: \url{https://github.com/Seetie-AI/memory_state}}

\begin{document}
\maketitle

\begin{abstract}

Agent memory systems typically pair a reasoning model with a separate,
smaller embedding model that indexes what the agent should remember. We
ask whether the reasoning model can supply its own retrieval keys.
Memory State prompts a generative LLM to compress a memory page into a
short key, then stores the hidden state at the final key token, taken
from a late-layer band near the output, as the retrieval vector. No
retriever is trained, distilled, or adapted.

On LongMemEval-S at round granularity, a compact two-view configuration
using Qwen3.5-9B-MLX-4bit reaches strict Recall@5 of 0.777 using hidden
states alone, at 8 KB per page, against 0.755 for a
Qwen3-Embedding-8B-4bit-DWQ baseline on the same 94 scored questions.
The margin is two questions on this subset, so we report it as
parity-level evidence and make no significance claim.

To test whether prompt semantics carry across benchmarks rather than
fitting one of them, we carry the three prompt identities selected on
LongMemEval over to PrefEval implicit-persona, a preference and persona
memory task, and evaluate them over all 1,000 examples. Per-view layer
and geometry settings were assigned on the PrefEval side, so this is
prompt-identity transfer rather than a frozen configuration. As pure
dense hidden-state retrieval with no lexical or embedding signal it
reaches Recall@5 of 0.299 against 0.281 for the same 8B embedding
baseline, while ranking last among the four K3 candidates evaluated
there. A model-size check under matched conditions gives 0.596 for
Qwen3.5-2B-MLX-4bit and 0.691 for Qwen3.5-9B-MLX-4bit, so in this
matched comparison the memory representation improves with the reasoning
model rather than being capped by a smaller retriever.

\end{abstract}

\section{Introduction}\label{introduction}

File-based agent memory usually begins with plain text and keyword
search. It is cheap, transparent, and easy to inspect. It also has a
blind spot: the agent must know what to search for. A memory that
matters to the current situation may share none of the words in the
current query. This is the unknown unknown problem.

Semantic retrieval addresses it, but it adds a second model. That
retriever is usually smaller than the reasoning model, and it was
trained for a retrieval objective that may not match the agent-memory
domain it is deployed in. If it misses, the more capable model never
sees the memory.

This work tests a third arrangement. The model that reads and reasons
over a memory also produces the vector that finds it later. If that
works, the memory index shares an inference path, a tokenizer, and a
semantic space with downstream reasoning, and it needs no separately
trained component.

The question came from an anomaly. While sweeping prompt suffixes on
PrefEval, a generic prompt asking the model for a single emoji reached
Recall@5 of about 0.285, next to 0.281 for an 8B embedding model on the
same 1,000 examples. That prompt was not the best in the sweep and it is
not part of any system reported here. We record it as the observation
that motivated the work: a reasoning model appears to carry a usable
memory representation in places we do not normally treat as an index.

We report three things. First, hidden-state keys reach
embedding-baseline scale on evidence retrieval. Second, a prompt
combination selected on one memory benchmark carries over to a different
memory task. Third, the representation improves with model size under
matched conditions.

\section{Related Work}\label{related-work}

\textbf{PromptEOL}~\cite{jiang2024} showed that a one-word prompt can turn the
hidden state of the final prompt token into a sentence embedding without
fine-tuning.

\textbf{MetaEOL}~\cite{lei2024} extended this with multiple meta-task prompt
views over the same text, and reported that the final layer is not
always the best choice, with a proportional layer-selection strategy
giving improvements.

\textbf{PromptReps}~\cite{zhuang2024} brought prompted LLM representations to
training-free full-corpus document retrieval, combining the last
layer\textquotesingle s last-token hidden state with next-token logits
into a dense and sparse hybrid, and reported stronger retrieval with
larger LLMs.

\textbf{All-but-the-Top}~\cite{mu2018} established that removing the common
mean and the top principal components from a set of embeddings exposes
semantic differences that a shared dominant direction otherwise masks.

Memory State does not claim the mechanism. Prompted final-token hidden
states, multiple prompt views, layer selection below the output layer,
and dominant-direction removal all come from this line of work. Our
contribution is the application: treating these representations as
persistent retrieval keys for long-term agent memory, designing
memory-specific prompt suffixes, and evaluating the result on two memory
benchmarks.

We evaluate on \textbf{LongMemEval}~\cite{wu2024} and
\textbf{PrefEval}~\cite{zhao2025}.

\section{Method}\label{method}

For each memory page and each query, the model is run over the text plus
a retrieval suffix that asks it to compress the text into a short key.
The retrieval vector is the hidden state at the final token of that
suffix.

The pattern is:

\begin{enumerate}
\def\labelenumi{\arabic{enumi}.}
\tightlist
\item
  choose prompt views that match the memory domain;
\item
  extract the suffix-final hidden state from a late-layer band near the
  output;
\item
  apply dominant-direction removal to strip shared prompt and corpus
  geometry;
\item
  combine views by concatenation or component-normalized vector
  averaging;
\item
  optionally add a lexical score as a shortlist rerank signal.
\end{enumerate}

Queries are encoded with the same recipe as memory pages. We tested the
asymmetric alternative, encoding memory and query with different
prompts, and it did not help: the best asymmetric configuration reached
Recall@5 of 0.755 against 0.766 for symmetric encoding of the same
prompt.

The reported configurations use hidden states as the dense
representation. Next-token logits were also collected, and a
PromptReps-style sparse path was reproduced, but it is not part of the
reported systems.

\subsection{Prompt views}\label{prompt-views}

Suffixes are written in Chinese and ask the model to emit a single
token. One of the PrefEval views, translated, reads:

\begin{verbatim}
Use one emoji to mark the emotion of the conversation above. The emoji is: "
\end{verbatim}

The trailing quotation mark is part of the suffix. It opens the
model\textquotesingle s answer, so the next decoded token is the key
itself. The other two views ask for one keyword the conversation should
recall and one token for the association it creates. The originals and
the full configuration are in the repository.

\subsection{Geometry}\label{geometry}

Dominant-direction removal is applied symmetrically to query and
candidate vectors with rank 15. Across 17 symmetric prompt variants,
symmetric removal won Recall@5 for 13, query-only removal for 3, and
plain centered cosine for 1.

\section{Experimental Setup}\label{experimental-setup}

Hidden-state extraction uses \texttt{Qwen3.5-9B-MLX-4bit}, with
\texttt{Qwen3.5-2B-MLX-4bit} for the model-size check. Embedding
baselines are \texttt{Qwen3-Embedding-8B-4bit-DWQ} and
\texttt{Qwen3-Embedding-0.6B-4bit-DWQ}, from the Qwen3 Embedding
series~\cite{zhang2025}. All models are local 4-bit conversions. Embedding scores are
never concatenated into the hidden vectors; they appear only as
baselines or as score-level signals.

LongMemEval-S is evaluated at round granularity on a 100-instance
subset, of which 94 questions score after the official abstention
filter. Recall@5 is strict: a question scores 1 only if every gold turn
appears in the top 5. One question moves the metric by about 1.1
percentage points. PrefEval implicit-persona is evaluated over all 1,000
examples.

Benchmark data is not redistributed with the repository. Sources and
download instructions are documented there.

\section{Results}\label{results}

\subsection{LongMemEval}\label{longmemeval}

\begin{table}[htbp]
\centering
\small
\caption{Retrieval results on LongMemEval-S.}
\label{tab:longmemeval}
\begin{tabularx}{\linewidth}{@{}Xrrr@{}}
\toprule
Method & R@5 & NDCG@5 & MRR \\
\midrule
BM25 & 0.585 & 0.547 & 0.559 \\
Qwen3-Embedding-8B-4bit-DWQ & 0.755 & 0.789 & 0.826 \\
Memory State, compact K2, hidden only, 8 KB/page & 0.777 & 0.788 &
0.820 \\
Memory State, K3 concat plus BM25 top 20, 24 KB/page & 0.777 & 0.822 &
0.851 \\
\bottomrule
\end{tabularx}
\begin{minipage}{\linewidth}
\vspace{3pt}\footnotesize\emph{Metric guide, for this and every table below. R@k (Recall@k)
measures how much relevant memory evidence appears within the first k
results; LongMemEval uses a strict version requiring every gold turn in
the top five. NDCG@5 rewards relevant evidence appearing nearer the top.
MRR measures how early the first relevant result appears. Higher is
better throughout.}
\end{minipage}
\end{table}

The compact two-view configuration matches the three-view hybrid on
strict recall while using no lexical signal and one third of the
storage. Its 0.022 margin over the 8B embedding baseline equals two
questions in this sample. At that resolution we make no significance
claim. This is parity-level evidence at embedding-baseline scale without
training an embedding model, not a decisive held-out win.

\subsection{Cross-task transfer to
PrefEval}\label{cross-task-transfer-to-prefeval}

The three prompt identities were selected on LongMemEval. They were then
evaluated on PrefEval as pure dense hidden-state retrieval over all
1,000 examples. Per-view layer and dominant-direction settings were
assigned when the PrefEval K3 candidates were constructed, so what
transfers is the prompt combination, not a frozen extraction
configuration.

\begin{table}[htbp]
\centering
\small
\caption{Cross-task retrieval results on PrefEval implicit-persona.}
\label{tab:prefeval}
\begin{tabularx}{\linewidth}{@{}Xrrrrr@{}}
\toprule
Method & R@1 & R@3 & R@5 & NDCG@5 & MRR \\
\midrule
BM25 & 0.035 & 0.074 & 0.094 & 0.066 & 0.067 \\
Qwen3-Embedding-8B-4bit-DWQ & 0.093 & 0.216 & 0.281 & 0.191 & 0.200 \\
Memory State, LongMem prompt transfer, hidden only & 0.093 & 0.212 &
0.299 & 0.197 & 0.188 \\
\bottomrule
\end{tabularx}
\end{table}

The transferred combination exceeds the 8B embedding baseline on strict
Recall@5 using no lexical signal and no embedding score. It is slightly
lower on R@3 and MRR. The LongMemEval shortlist and fusion weights were
not carried over.

Two things make this the most informative PrefEval row we have. It was
not selected to win here: of the four K3 candidates evaluated on
PrefEval it placed last, below combinations whose prompts were chosen on
PrefEval itself. And its prompt identities were fixed by a different
benchmark. So it is the weakest candidate in its own sweep and still
clears the 8B embedding baseline on strict recall. This supports the
claim that a prompt combination can cross memory domains, not that every
low-level parameter is universal.

A separate configuration tuned directly on PrefEval, fusing the
hidden-state score with an embedding score and BM25, reaches Recall@5 of
0.332. Because its prompts and fusion weights were selected on the same
data it is reported on, we treat it as a domain-tuned exploratory result
and not as the headline.

\subsection{Model size}\label{model-size}

Both models were run with the same prompt suffix, the same centered
transform, the same 4-bit precision, and the same 100 instances.

\begin{table}[htbp]
\centering
\small
\caption{Matched model-size comparison.}
\label{tab:model-size}
\begin{tabularx}{0.8\linewidth}{@{}Xrrr@{}}
\toprule
Model & Layer & R@5 & NDCG@5 \\
\midrule
Qwen3.5-2B-MLX-4bit & 22 & 0.596 & 0.547 \\
Qwen3.5-9B-MLX-4bit & 30 & 0.691 & 0.689 \\
\bottomrule
\end{tabularx}
\end{table}

The larger model gains 9.5 percentage points on Recall@5. On a smaller
30-instance tuning slice both models scored 0.833, so the gap appears
only at the larger sample. Each model uses its own best late layer
rather than a fixed absolute index. The direction is consistent with
earlier retrieval results reported for prompted representations at
larger model scale~\cite{zhuang2024}.

\section{Analysis and Negative
Results}\label{analysis-and-negative-results}

\textbf{There is no universal best layer.} Across the prompt sweep,
layers 29, 30, and 31 each won for 6, 6, and 5 variants respectively.
The useful region is a late-layer band whose exact position depends on
prompt semantics, not a single layer.

\textbf{Prompt wording matters more than expected.} Changing one word in
an otherwise identical suffix, from a neutral term for the other party
to the product-native term for the user, moved Recall@5 by 0.043. The
model appears better aligned to product vocabulary than to neutral
phrasing.

\textbf{Prompt language is not neutral.} In the sweep, the best cell for
an English variant of one suffix reached Recall@5 of 0.713 while the
best cell for its Chinese original reached 0.755, on English memory
content. The two best cells do not share a layer and transform, so this
is a sweep observation rather than a single-variable comparison.
Cross-language prompting worked here; we make no general multilingual
claim.

\textbf{Not every retrieval-shaped prompt works.} An answer-strategy
prompt fell to Recall@5 of 0.574.

\textbf{Lexical fusion helps only as a small vote inside a shortlist.}
Fusing BM25 inside the vector top 20 or top 50 beat full-corpus fusion
for ranking, and weighting BM25 heavily collapsed performance across
configurations, because turn-level lexical matching is noisy.

\textbf{Late interaction did not justify its cost.} Max-sim style
scoring over multiple stored vectors did not beat the simpler
concatenation and averaging candidates.

\textbf{Agreement filtering is a diagnostic, not a default.} Gold
candidates were much more likely than non-gold candidates to appear in
multiple prompt rankings, but sorting by agreement never beat sorting by
score.

\textbf{A stronger reranker buys ranking, not recall.} Blending an 8B
embedding score into the hybrid raised NDCG@5 to 0.834 and MRR to 0.888
without improving Recall@5, while reintroducing the separate model
dependency the method exists to avoid. We keep it as a ceiling, not a
default.

\textbf{Small subsets misled us.} An early prompt that led a 40-instance
slice fell to mid-pack on the full 100 instances. Oracle union
diagnostics, which use gold labels to find complementary prompt pairs,
were likewise treated as hypothesis generators rather than estimates.

\section{Deployment Considerations}\label{deployment-considerations}

The vector comes from the final token of a prompted key-generation pass.
In an online setting the reasoning model has already processed the
conversation and its context is in the KV cache, so capturing a memory
key means decoding one additional token per prompt view and keeping its
hidden state. Memory capture then happens inside the inference path
rather than as a second encoding job by another model. We did not
measure end-to-end latency or cost, so this is a deployment argument,
not a measured result.

Storage is straightforward. Qwen3.5-9B has a hidden size of 4096, so one
view in bf16 costs 8 KB per page and three views cost 24 KB. Around
20,000 memory pages need roughly 500 MB before index overhead. The
compact configuration in Section 5.1 reaches the same strict recall as
the three-view hybrid at 8 KB.

\section{Limitations}\label{limitations}

The LongMemEval numbers rest on 94 scored questions, where one question
moves the metric by about 1.1 percentage points. Differences of one or
two questions between the reported systems and the embedding baselines
should be read as parity, not as ordering.

Prompt and fusion choices were selected on explored benchmark subsets.
The LongMemEval results are therefore exploratory rather than a clean
held-out evaluation. The PrefEval result in Section 5.2 is the closest
we come to an independent estimate, because its prompt identities were
fixed by a different benchmark and it was not the configuration selected
to perform best there. Its per-view layer and geometry settings were
still assigned on the PrefEval side, so it is not a held-out evaluation
either.

The geometry is tied to \texttt{Qwen3.5-9B-MLX-4bit} and may change
under another model, another quantization, another language, or another
prompt template. Both models tested come from one family, so we cannot
say whether the effect holds across model families.

The work supports feasibility at embedding-baseline scale. It does not
establish a general win over embedding models.

\section{Reproducibility}\label{reproducibility}

Code, experiment logs, the earlier internal results overview, and the
result files for all reported numbers are public at
\url{https://github.com/Seetie-AI/memory_state}.

Benchmark datasets and model weights are not redistributed. The
repository documents the original sources for LongMemEval and PrefEval,
the download commands, and the model identifiers, so that users obtain
them under the original terms.

\section{Conclusion}\label{conclusion}

A generative model, in the course of reading a page, leaves behind a
representation that can be used to find that page again. On two
different memory benchmarks this representation performed at the scale
of embedding models built for the job, without training a retriever, and
it improved with model size under matched conditions.

The claim is not that hidden states beat embedding models. It is that an
agent\textquotesingle s memory index may not require a second model at
all, and that the memory an agent needs may already be present in the
act of understanding, waiting to be read out rather than rebuilt.

\begin{thebibliography}{9}

\bibitem{jiang2024}
Jiang et al.
\newblock Scaling Sentence Embeddings with Large Language Models.
\newblock \emph{Findings of EMNLP}, 2024.
\newblock \url{https://aclanthology.org/2024.findings-emnlp.181/}

\bibitem{lei2024}
Lei et al.
\newblock Meta-Task Prompting Elicits Embeddings from Large Language Models.
\newblock \emph{ACL}, 2024.
\newblock \url{https://aclanthology.org/2024.acl-long.546/}

\bibitem{zhuang2024}
Zhuang et al.
\newblock PromptReps: Prompting Large Language Models to Generate Dense and
Sparse Representations for Zero-Shot Document Retrieval.
\newblock \emph{EMNLP}, 2024.
\newblock \url{https://aclanthology.org/2024.emnlp-main.250/}

\bibitem{mu2018}
Mu and Viswanath.
\newblock All-but-the-Top: Simple and Effective Postprocessing for Word
Representations.
\newblock \emph{ICLR}, 2018.
\newblock \url{https://openreview.net/forum?id=HkuGJ3kCb}

\bibitem{wu2024}
Wu et al.
\newblock LongMemEval: Benchmarking Chat Assistants on Long-Term Interactive
Memory.
\newblock 2024. \url{https://arxiv.org/abs/2410.10813}

\bibitem{zhao2025}
Zhao, Hong, Liu, Hazarika, and Lin.
\newblock Do LLMs Recognize Your Preferences? Evaluating Personalized
Preference Following in LLMs.
\newblock 2025. \url{https://arxiv.org/abs/2502.09597}

\bibitem{zhang2025}
Zhang et al.
\newblock Qwen3 Embedding: Advancing Text Embedding and Reranking Through
Foundation Models.
\newblock 2025. \url{https://arxiv.org/abs/2506.05176}

\end{thebibliography}

\end{document}
